\chapter{Exercises}
\label{ch:exercises}

\textit{Six hands-on exercises that progressively deepen your understanding
of \sysname{}.  Each exercise states a learning objective, difficulty
level, estimated time, the files you will work in, step-by-step
instructions, and acceptance criteria.  Hints are provided at the end of
each exercise for when you get stuck.}

% ─────────────────────────────────────────────────────────────────────
\section{Ontology Expansion}
\label{sec:ex-ontology-expansion}

\noindent\textbf{Difficulty:} Easy \quad
\textbf{Time:} $\sim$2 hours \quad
\textbf{Learning objective:} Understand canonical mappings and the
synthetic benchmark.

\medskip\noindent\textbf{Files.}

\begin{itemize}
    \item \codefile{apps/api/models/ontology.py}
    \item \codefile{apps/api/evaluation/synthetic\_benchmark.py}
\end{itemize}

\medskip\noindent\textbf{Instructions.}

\begin{enumerate}
    \item Pick a technical domain you know well (e.g.\ game development,
          embedded systems, bioinformatics).
    \item Add 20 new canonical mappings to the \ic{CANONICAL\_NAMES}
          dictionary in \codefile{ontology.py}.  For each canonical name,
          add at least two surface-form aliases (e.g.\ \ic{"unreal"}
          $\rightarrow$ \ic{"Unreal Engine"},
          \ic{"ue5"} $\rightarrow$ \ic{"Unreal Engine"}).
    \item For any canonical name that has well-known abbreviations, add
          them to \ic{KNOWN\_ABBREVIATIONS} in
          \codefile{synthetic\_benchmark.py}.
    \item Run the benchmark and confirm all new mappings pass.
\end{enumerate}

\begin{tryit}
\textbf{Acceptance criteria:}
\begin{itemize}
    \item 20 new entries in \ic{CANONICAL\_NAMES}, each with a lowercase
          key and properly cased canonical value.
    \item \ic{python -m evaluation.synthetic\_benchmark} runs to
          completion with no failures.
    \item No existing mappings are broken (the benchmark tests all of
          them).
\end{itemize}
\end{tryit}

\medskip\noindent\textbf{Hints.}

\begin{itemize}
    \item Look at the existing category sections (\ic{DATABASES},
          \ic{AI/ML FRAMEWORKS}, etc.) for formatting conventions.
    \item Remember: keys must be lowercase.  The resolver calls
          \ic{name.lower()} before lookup.
    \item If you are unsure whether a name already exists, search the file
          for the canonical form first.
\end{itemize}


% ─────────────────────────────────────────────────────────────────────
\section{Graph Layout Enhancement}
\label{sec:ex-graph-layout}

\noindent\textbf{Difficulty:} Easy--Medium \quad
\textbf{Time:} $\sim$3 hours \quad
\textbf{Learning objective:} Work with the React Flow graph renderer and
understand how node positions are computed.

\medskip\noindent\textbf{Files.}

\begin{itemize}
    \item \codefile{apps/web/} --- graph page component and layout
          utilities.
\end{itemize}

\medskip\noindent\textbf{Instructions.}

\begin{enumerate}
    \item Study the existing graph layout code.  The default layout
          arranges nodes using a force-directed or hierarchical algorithm.
    \item Implement a \textbf{circular layout} option that places nodes in
          a circle with the most-connected node at the top.
    \item Add a layout toggle (dropdown or button group) to the graph page
          so users can switch between the existing layout and your new
          circular layout.
    \item Ensure edges render correctly with the new positions.
\end{enumerate}

\begin{tryit}
\textbf{Acceptance criteria:}
\begin{itemize}
    \item A ``Circular'' option appears in the graph page layout controls.
    \item Selecting it repositions all nodes in a circle.
    \item The most-connected node (highest degree) appears at the 12
          o'clock position.
    \item Edges remain visible and do not overlap excessively.
    \item Switching back to the default layout restores the original
          positions.
    \item \ic{pnpm typecheck} passes with no errors.
\end{itemize}
\end{tryit}

\medskip\noindent\textbf{Hints.}

\begin{itemize}
    \item The circular position for node $i$ out of $N$ total is:
          $x = r \cos(2\pi i / N)$, $y = r \sin(2\pi i / N)$.
    \item Sort nodes by degree (number of connections) descending before
          assigning positions so the highest-degree node gets index~0.
    \item Use \ic{cn()} from \ic{@/lib/utils} for conditional class names
          on the toggle buttons.
\end{itemize}


% ─────────────────────────────────────────────────────────────────────
\section{Search Entity Type Filter}
\label{sec:ex-search-filter}

\noindent\textbf{Difficulty:} Medium \quad
\textbf{Time:} $\sim$3 hours \quad
\textbf{Learning objective:} Add a feature that spans both backend and
frontend, touching an API endpoint and a UI control.

\medskip\noindent\textbf{Files.}

\begin{itemize}
    \item \codefile{apps/api/routers/search.py} --- add query parameter.
    \item \codefile{apps/api/services/} --- filter logic in the search
          service.
    \item \codefile{apps/web/} --- search page component.
\end{itemize}

\medskip\noindent\textbf{Instructions.}

\begin{enumerate}
    \item Add an optional \ic{entity\_type} query parameter to the search
          endpoint in \codefile{routers/search.py}.  Accept any value from
          the \ic{EntityType} enum (e.g.\ \ic{technology}, \ic{concept}).
    \item Update the search service to filter results by entity type when
          the parameter is present.  When absent, return all types (current
          behavior).
    \item On the frontend search page, add a dropdown (or segmented
          control) that lists all entity types plus an ``All'' option.
    \item Wire the dropdown selection into the search API call as a query
          parameter.
    \item Test with several entity types and verify the result count
          changes.
\end{enumerate}

\begin{tryit}
\textbf{Acceptance criteria:}
\begin{itemize}
    \item \ic{GET /api/search/?q=react\&entity\_type=technology} returns
          only technology entities.
    \item \ic{GET /api/search/?q=react} (no filter) returns all types.
    \item The frontend dropdown updates search results in real time.
    \item Invalid entity types return a 422 validation error.
    \item \ic{.venv/bin/pytest tests/ -v -k search} passes.
    \item \ic{pnpm typecheck} passes.
\end{itemize}
\end{tryit}

\medskip\noindent\textbf{Hints.}

\begin{itemize}
    \item Use \ic{Optional[EntityType] = None} as the FastAPI parameter
          type for automatic enum validation.
    \item The Neo4j query likely uses a \ic{WHERE} clause on node labels
          or a \ic{type} property---add a conditional filter there.
    \item On the frontend, use React Query's \ic{queryKey} to include the
          entity type so cache invalidation works correctly when the filter
          changes.
\end{itemize}


% ─────────────────────────────────────────────────────────────────────
\section{New Evaluation Baseline}
\label{sec:ex-new-baseline}

\noindent\textbf{Difficulty:} Medium \quad
\textbf{Time:} $\sim$4 hours \quad
\textbf{Learning objective:} Implement a comparison baseline for entity
resolution using traditional NLP techniques.

\medskip\noindent\textbf{Files.}

\begin{itemize}
    \item \codefile{apps/api/evaluation/tfidf\_baseline.py} --- new file.
    \item \codefile{apps/api/evaluation/data/v5/} --- output directory.
\end{itemize}

\medskip\noindent\textbf{Instructions.}

\begin{enumerate}
    \item Create \codefile{evaluation/tfidf\_baseline.py}.
    \item Load the synthetic benchmark data (the CSV produced by
          \codefile{synthetic\_benchmark.py}).
    \item For each surface-form variant, compute its TF-IDF vector using
          scikit-learn's \ic{TfidfVectorizer} with character n-grams
          (\ic{analyzer="char\_wb"}, \ic{ngram\_range=(2,4)}).
    \item Build a TF-IDF matrix for all canonical names.
    \item For each test variant, find the nearest canonical name by cosine
          similarity and predict that as the resolution.
    \item Compute precision, recall, and F1 against the ground-truth
          canonical labels.
    \item Write results to
          \codefile{evaluation/data/v5/tfidf\_baseline\_results.json}.
\end{enumerate}

\begin{tryit}
\textbf{Acceptance criteria:}
\begin{itemize}
    \item The script runs end-to-end:
          \ic{python -m evaluation.tfidf\_baseline}.
    \item Output JSON contains \ic{precision}, \ic{recall}, and \ic{f1}.
    \item Results are computed on the held-out test split only (60
          conversations, not all 200).
    \item The script prints a summary comparing TF-IDF F1 against
          \sysname{}'s pipeline F1 (0.979 from B-cubed evaluation).
\end{itemize}
\end{tryit}

\medskip\noindent\textbf{Hints.}

\begin{itemize}
    \item Install scikit-learn: \ic{.venv/bin/pip install scikit-learn}.
    \item Character n-grams work better than word n-grams for entity name
          matching because many tech names are single tokens.
    \item Use \ic{cosine\_similarity} from
          \ic{sklearn.metrics.pairwise} for efficient batch comparison.
    \item The existing benchmark CSV has a \ic{split} column
          (\ic{train}/\ic{test}) to identify held-out data.
\end{itemize}


% ─────────────────────────────────────────────────────────────────────
\section{GraphRAG Depth Experiment}
\label{sec:ex-graphrag-depth}

\noindent\textbf{Difficulty:} Medium--Hard \quad
\textbf{Time:} $\sim$4 hours \quad
\textbf{Learning objective:} Understand how graph traversal depth affects
retrieval quality in the GraphRAG pipeline.

\medskip\noindent\textbf{Files.}

\begin{itemize}
    \item \codefile{apps/api/services/graph\_rag.py} --- retrieval depth
          parameter.
    \item \codefile{apps/api/evaluation/} --- new experiment script.
\end{itemize}

\medskip\noindent\textbf{Instructions.}

\begin{enumerate}
    \item Study \codefile{services/graph\_rag.py} to understand how the
          retrieval pipeline expands from seed entities through graph
          neighbors.
    \item Create an experiment script that runs the same set of 10 queries
          at depths 1, 2, and 3.
    \item For each depth, record: number of nodes retrieved, number of
          unique entity types, and the context string length.
    \item Manually judge relevance of the top-5 retrieved nodes for each
          query at each depth (use a simple 0/1/2 scale: irrelevant,
          partially relevant, highly relevant).
    \item Compute average relevance scores per depth and plot the
          trade-off between depth and relevance.
\end{enumerate}

\begin{tryit}
\textbf{Acceptance criteria:}
\begin{itemize}
    \item The experiment script produces a JSON file with per-query,
          per-depth results.
    \item At least 10 queries are tested at all three depths.
    \item A summary table shows: depth, avg nodes retrieved, avg
          relevance, avg context length.
    \item The script includes a brief analysis comment interpreting whether
          deeper traversal helps or hurts retrieval quality.
\end{itemize}
\end{tryit}

\medskip\noindent\textbf{Hints.}

\begin{itemize}
    \item The \ic{retrieve\_context()} method in \ic{GraphRAGService}
          returns \ic{(subgraph, context\_str, seed\_ids)}.  The subgraph
          contains the retrieved nodes.
    \item Depth~1 retrieves direct neighbors; depth~2 retrieves
          neighbors-of-neighbors.
    \item You will likely find that depth~2 is the sweet spot---depth~3
          tends to bring in too much noise for most queries.
    \item The existing \codefile{evaluation/test\_graphrag.py} has example
          queries you can reuse.
\end{itemize}


% ─────────────────────────────────────────────────────────────────────
\section{New MCP Tool}
\label{sec:ex-mcp-tool}

\noindent\textbf{Difficulty:} Hard \quad
\textbf{Time:} $\sim$4 hours \quad
\textbf{Learning objective:} Extend the Model Context Protocol server
with a new tool and write automated test cases.

\medskip\noindent\textbf{Files.}

\begin{itemize}
    \item \codefile{apps/mcp/} --- MCP server tool definitions.
    \item \codefile{apps/api/evaluation/test\_mcp.py} --- MCP test suite.
\end{itemize}

\medskip\noindent\textbf{Instructions.}

\begin{enumerate}
    \item Design a \ic{continuum\_timeline} MCP tool that accepts an
          entity name and returns a chronological list of decisions
          involving that entity.
    \item The tool should query the Neo4j graph for all
          \ic{DecisionTrace} nodes connected to the entity via
          \ic{INVOLVES} relationships, sorted by timestamp.
    \item Each entry in the returned list should include: decision
          summary, timestamp, confidence score, and any
          \ic{SUPERSEDES}/\ic{CONTRADICTS} relationships to other
          decisions.
    \item Implement the tool in the MCP server following the patterns of
          existing tools.
    \item Write at least 5 test cases covering:
          \begin{itemize}
              \item Entity with no decisions (empty result).
              \item Entity with one decision (single entry).
              \item Correct chronological ordering (oldest first).
              \item Entity with a \ic{SUPERSEDES} chain.
              \item Entity with a \ic{CONTRADICTS} pair.
          \end{itemize}
\end{enumerate}

\begin{tryit}
\textbf{Acceptance criteria:}
\begin{itemize}
    \item The \ic{continuum\_timeline} tool appears in the MCP server's
          tool list.
    \item Calling it with \ic{"PostgreSQL"} returns a chronological list
          of decisions.
    \item The response includes \ic{SUPERSEDES}/\ic{CONTRADICTS} metadata
          when present.
    \item All 5 test cases pass when running the MCP test suite.
    \item The tool handles non-existent entity names gracefully (returns
          empty list, not an error).
\end{itemize}
\end{tryit}

\medskip\noindent\textbf{Hints.}

\begin{itemize}
    \item Study the existing MCP tools in \codefile{apps/mcp/} for the
          registration pattern and response format.
    \item The Cypher query will need to \ic{MATCH} the entity, traverse
          \ic{INVOLVES} relationships to decision nodes, and optionally
          \ic{MATCH} inter-decision relationships.
    \item Use \ic{ORDER BY d.created\_at ASC} (or the equivalent timestamp
          property) for chronological sorting.
    \item The existing \codefile{evaluation/test\_mcp.py} has 31 test
          cases that serve as a template for test structure.
\end{itemize}
